A planar Geyer and Herr reflex walking model in SCONE was given right drop foot by weakening tibialis anterior to 50, 25 and 10\% of its strength. With its healthy controller it fell at every level below 96\%; after re-optimization of an asymmetric controller it walked at all three with a steppage gait, a plantarflexed foot and a 5 to 11\% higher cost of transport. For 25\% strength an ankle exoskeleton was built end to end: a torque requirement from the lost muscle moment, a maxon EC~45 flat motor with gear ratio 75 and a \SI{1081}{N.m/rad} series spring sized for energy per stride, a PID plus feedforward torque controller with \SI{29}{Hz} bandwidth reduced to a delay-and-lag model, and a causal heel-strike and toe-off detector on a simulated \SI{100}{Hz} shank gyroscope with noise, bias and delay. The detector found 99\% of events on simulated drop-foot gait and, unchanged, on real gyroscopes of 20 adults. A hypothesis committed before the comparison, that a phase-based active AFO restores toe clearance without lowering push-off power while the best passive AFO lowers it by more than 15\%, was rejected over three optimization seeds: the passive AFO's spring returned energy at push-off and raised the total push-off power. Neither device could be worn without re-optimizing the reflexes. After re-optimization the active AFO did not change toe clearance, which the adapted model already over-compensated, but let the model drop most of its steppage and lowered the cost of transport by 5\%.
